\documentclass{article}
\PassOptionsToPackage{hyphens}{url}
\usepackage{iclr2025_conference,times}
\usepackage{hyperref}
\usepackage{xurl}

\hypersetup{hidelinks,
  pdftitle={Is the Gap an Interface Gap? An Agent Tool Interfaces for Data Catalogs},
  pdfauthor={Alex Feng}}

% CS329Z project proposal, revision 16: two-page prose version (body only; references and AI disclosure excluded).

\title{[Midway Report] Is the Gap an Interface Gap?\\ An Agent Tool Interfaces for Data Catalogs}
\author{Alex Feng}
\iclrfinalcopy

\begin{document}
\maketitle
\vspace{-0.3in}
\lhead{CS329Z: Engineering AI Agents -- Project Proposal}

\section{Problem Statement}



\section{Methodology}



\section{Results}



\section{Assets}



\section{Timeline}



\clearpage
\bibliography{references}
\bibliographystyle{iclr2025_conference}

\appendix
\section{AI-Use Disclosure Statement}
In this work, I used generative AI tools to refine the
hypothesis, to brainstorm task formulations, and to
give feedback on the research methodology. The remaining required-disclosure tasks
(generating or cleaning datasets, mathematical claims or proofs,
implementing methods, translation, qualitative data analysis, and
interpreting results) are not applicable given that this is just a proposal rather than the final study. Additionally, I used generative AI
tools for brainstorming, identifying gaps in prior work, searching for
and summarizing relevant literature, and
providing editing suggestions for readability. I have reviewed
all AI-assisted work, including the literature summaries, for accuracy, and I take
responsibility for the final content of this work, including text,
claims, or artifacts produced with the aid of generative AI.

<this needs to be updated^^>
\end{document}